\chapter{Lecture 15}

\hfill {\footnotesize \emph{Date: 01.05.2026}}

\section{Integral Dependence}

Let \( R \) be a ring, and let  \( S \) be a subring of  \( R \) (so \( 1 \in S\)). An element \( x \in R \) is said to be  \emph{integral over }\( S \) if  \( x \) is a root of a monic polynomial with coefficients from  \( S \)---so that means  \( x \) satisfies an equation of the form
 \[
   x^n + a_1x^{n-1} + \cdots + a_n = 0, \quad a_i \in S, \forall i \in \{1, \ldots, n\}.
\]
\begin{remark}
  Every element \( s \in S \) is integral over \( S \) because 
  \[
  s^2 + (-s)s + 0 = 0
  \]
\end{remark}
\begin{example}\label{ex:integralclosureZ}
  Suppose \(R = \mathbb{Q}\) and \(S = \mathbb{Z}\). We want to find rational numbers in \(\mathbb{Q}\) that are integral over \(\mathbb{Z}\).

  Let \(\frac{r}{s} \in \mathbb{Q}\) be in lowest terms, i.e., \(\gcd(r,s)=1\). If \(\frac{r}{s}\) is integral over \(\mathbb{Z}\), then it satisfies a monic polynomial equation with integer coefficients:
\[
\left(\frac{r}{s}\right)^{\!n} + a_1\left(\frac{r}{s}\right)^{\!n-1} + \dots + a_n = 0, a_i \in \mathbb{Z}.
\]
Multiplying both sides by \(s^n\) yields:
\begin{gather*}
  r^n + a_1 r^{n-1} s + a_2 r^{n-2} s^2 + \dots + a_n s^n = 0\\
  \implies r^n = -s\bigl(a_1 r^{n-1} + a_2 r^{n-2} s + \dots + a_n s^{n-1}\bigr).
\end{gather*}
Hence \(s\) divides \(r^n\).  Because \(\gcd(r,s)=1\), we also have \(\gcd(r^n,s)=1\).  The only way \(s \mid r^n\) and \(\gcd(r^n,s)=1\) is that \(s = \pm 1\). Therefore \(\frac{r}{s} \in \mathbb{Z}\). This shows that every rational number integral over \(\mathbb{Z}\) is an integer.
\end{example}
We will now state and prove some equivalent conditions of being an integral element. But before that, a definition: 
\begin{definition}[Faithful \( R \)-module]
  An \( R \)-module  \( N \) is called  \emph{faithful} if for all \( R \)-module  \( M \), \( M \oplus_R N = 0 \) implies  \( M = 0 \).
\end{definition}
\todo{Did not understand the defintion}
\begin{proposition}\label{Prop:EquivIntegral}
  The following are equivalent:
  \begin{enumerate}[label=\roman*)]
    \item \( x \in R \) is integral over  \( S \).
    \item  \( S[x] \) is finitely generated as an  \( S \)-module.
    \item \( S[x] \) is contained in a subring  \( T \) of  \( R \) such that  \( T \) is finitely generated.
    \item There exists a faithful  \( S[x] \)-module  \( M \) which is finitely generated as an  \( S \)-module.
  \end{enumerate}
  Here \( S[x] \) means the set \( \{a_0x^n + a_1x^{n-1} + \cdots + a_n \mid a_i \in S, n \in \mathbb{N}\} \), not the polynomial ring over \( S \).
\end{proposition}
\begin{proof} The strategy is as follows:
\[
  (\text{i}) \underbrace{\implies}_{\textit{Step 01}} (\text{ii}) \underbrace{\implies}_{\textit{Step 02}} (\text{iii}) \underbrace{\implies}_{\textit{Step 03}}  (\text{iv}) \underbrace{\implies}_{\textit{Step 04}} (\text{i})  
\] 
  \emph{Step 01.} If \(x\) is integral over \(S\), then there exist \(a_1,\dots,a_n \in S\) such that
  \begin{gather*}
x^n + a_1 x^{n-1} + \dots + a_n = 0\tag{1}\\
\implies x^n = (-a_1) x^{n-1} + \dots + (-a_n).
\end{gather*}

Hence \(x^n\) belongs to the submodule of \(S[x]\) generated by \(1, x, \dots, x^{n-1}\).

Multiplying both sides of (1) by \(x^r\) (for \(r \ge 0\)) gives
\[
x^{n+r} = (-a_1) x^{n+r-1} + \dots + (-a_n) x^r.
\]

Now for example, with \(r=1\):
\[
x^{n+1} = (-a_1) x^n + (-a_2) x^{n-1} + \dots + (-a_n) x.
\]

Substituting the expression for \(x^n\) from (1) gives
\[
x^{n+1} \in \langle x^{n-1}, x^{n-2}, \dots, 1 \rangle.
\]

Therefore, we can inductively show that \(x^{n+r} \in \langle x^{n-1}, \dots, 1 \rangle\) for every \(r \ge 0\).

Now,
\[
S[x] = \langle 1, x, x^2, \dots \rangle = \langle 1, x, \dots, x^{n-1}, x^n, x^{n+1}, \dots \rangle.
\]

Since every power \(x^{n+r}\) lies in the \(S\)-span of \(1, x, \dots, x^{n-1}\), we conclude
\[
S[x] = \langle 1, x, \dots, x^{n-1} \rangle.
\]

Thus \(S[x]\) is a finitely generated \(S\)-module, generated by \(1, x, \dots, x^{n-1}\).

\noindent
\emph{Step 02.} Assuming that \( S[x] \) is finitely generated as an  \( S \)-module, take  \( T = S[x] \). 

\noindent
Step 03. 
\end{proof}
\begin{lemma} \label{lem:transModuleGen}
  If \( N \) is finitely generated as a  \( B \)-module and that  \( B \) is finitely generated as an  \( A \)-module. Then  \( N \) is finitely generated as an  \( A \)-module.
\end{lemma}
\begin{proof}
  \todo{To be written...}
\end{proof}
\begin{corollary}\label{cor:equivIntegralMulVersion}
  Let \( x_i \) \((1 \le i \le n) \) be elements of  \( R \), each integral over  \( S \). Then the ring  \( S[x_1, \ldots, x_n] \) is a finitely generated \( S \)-module.
\end{corollary}
\begin{proof}
  We use induction on \( n \).  \( n = 1 \) case was proven in (ii) of \autoref{Prop:EquivIntegral}. Assume the statement is true for \( n-1 \) where \( n > 1 \), let  \( S_r = S[x_1, \ldots, x_r] \); then by the inductive hypothesis \( S_{n-1} \) is a finitely generated  \( S \)-module. Since \( x_n \) is integral over  \( S \), it satisfies a monic polynomial with coefficients in  \( S \subseteq S[x_1, \ldots, x_{n-1}] \). Hence, it is integral over \( S_{n-1} \) as well. Hence, by the case \( n=1 \),  \( S_n = S_{n-1}[x_n] \) is a finitely generated as  \( S_{n-1} \)-module. Using \autoref{lem:transModuleGen}, \( S_n = [x_1, \ldots, x_n] \) is a finitely generated \( S \)-module.
\end{proof}
\begin{corollary}\label{cor:IntClSubring}
  Suppose \( R \) is a  ring and \( S \) is a subring of  \( R \). Then let  \( C \) be the set of all elements in  \( R \) which are integral over  \( S \). Then is  \( C \) is a subring of  \( R \) containing  \( S \).
\end{corollary}
\begin{proof}
  \( S \subseteq C \) because every element in \( S \) is an integral over  \( S \).

  Since \( S \subseteq C \), we get \( 0 \in C \) for free. If  \( x, y \in C \), then we want to show that  \( x - y \in C \) and  \( xy \in C\). Since \( x, y \in C \)---set of all ``integrals'' over  \( S \), using \autoref{cor:equivIntegralMulVersion}, \( S[x, y] \) would be finitely generated as an  \( S \)-module.

  We have \( S[x-y], S[xy]\) is contained in a subring  \( T = S[x,y] \) of  \( R \) such that  \( T \) is finitely generated as an  \( S \)-module. Using (iii) of \autoref{Prop:EquivIntegral}, \( x-y \), and  \( xy \) are integral over S.  
\end{proof}
In \autoref{cor:IntClSubring}, the set \( C \) of elements of \( R \) that are integral over \( S \) is a subring of  \( R \). This subring of  \( R \) is called  \emph{integral closure} of \( S \) in  \( R \). If the integral closure of \( S \) in  \( R \) is  \( S \) itself, then  \( S \) is called \emph{integrally closed}. In \autoref{ex:integralclosureZ}, we could see \( \mathbb{Z} \) is integrally closed in \( \mathbb{Q} \). Moreover, if the integral closure of  \( S \) in  \( R \) is  \( R \)  itself, then \( R \) is said to be  \emph{integral over} \( S \).
\begin{corollary}[Transitivity of integral dependence]\label{cor:IntDepTrans}
  Suppose \( A \subseteq B \subseteq C \) are rings, where \( B \) is integral over  \( A \) and  \( C \) is integral over  \( B \). Then  \( C \) is also integral over  \( A \). 
\end{corollary}
Here, \( B \) is integral over  \( A \) (respectively, \( C \) is integral over  \( B \)) means each \( b \in B \) is integral over \( A \) (respectively, each \( c \in C \) is integral over  \( B \)).
\begin{proof}[Proof of corollary]
  Let \(x \in C\). Since \(C\) is integral over \(B\), there exist \(b_1,\dots,b_n \in B\) such that
\[
x^n + b_1 x^{n-1} + b_2 x^{n-2} + \dots + b_n = 0. \tag{1}
\]
Now, \(B\) is integral over \(A\). Therefore the ring \(A[b_1,\dots,b_n]\) is a finitely generated \(A\)-module. Denote \(B' = A[b_1,\dots,b_n]\), then \(b_1,\dots,b_n \in B'\), and \( x^n + b_1x^{n-1} + \ldots + b_n = 0 \), so \(x\) is also integral over \(B'\) as well, hence \(B'[x]\) is finitely generated as a \(B'\)-module. Since \(B'\) is finitely generated as an \(A\)-module, it follows that \(B'[x]\) is a finitely generated \(A\)-module. Now
\[
B'[x] = A[b_1,\dots,b_n][x] = A[b_1,\dots,b_n,x].
\]

is finitely generated as an \( A \)-module. 

Consider the \(A\)-module homomorphism
\[
\phi : A[b_1,\dots,b_n,x] \to A[b_1,\dots,b_n,x], \phi(t) = t \cdot x.
\]

Then 
\[
  \phi(A[b_1, \ldots, b_n, x]) \subseteq A[b_1, \ldots, b_n, x] \subseteq A \cdot A[b_1, \ldots, b_n, x].
\] 

\( A \) is an ideal in itself as a ring; by \autoref{prop:fingeneratedmodpoly},  there exists \( a_1, \ldots, a_k \) such that
  \begin{align*}
    & \phi^k + a_1\phi^{k-1} + \ldots + a_k = 0\\
    \implies & (\phi^k + a_1\phi^{k-1} + \cdots + a_k)(1) = 0\\
    \therefore~&x^k + a_1x^{k-1} + \cdots + a_k = 0 \text{ for some }  a_1, \ldots, a_k \in A
  \end{align*} 
which implies \( x \) is integral over  \( A \), and so  \( C \) is integral over  \( A \).
\end{proof}
%   \begin{remark}
%   If \(x_1,\dots,x_n \in R\) are each integral over a subring \(S \subset R\), then the ring \(S[x_1,\dots,x_n]\) is a finitely generated \(S\)-module. The argument is inductive:
%     \begin{enumerate}
%       \item \(x_1\) integral over \(S\) if and only if \(S[x_1]\) is finitely generated as an \(S\)-module.
%       \item Since \(x_2\) is integral over \(S\), it is also integral over \(S[x_1]\). Therefore \((S[x_1])[x_2] = S[x_1,x_2]\) is finitely generated as an \(S[x_1]\)-module.
%       \item As \(S[x_1]\) is a finitely generated \(S\)-module, \(S[x_1,x_2]\) is finitely generated as an \(S\)-module.
%       \item Continue inductively: \(S[x_1,\dots,x_n]\) is a finitely generated \(S\)-module.
%     \end{enumerate}
%   \end{remark}
